\documentclass[12pt, a4paper]{article}
\usepackage[spanish]{babel}
\usepackage[utf8]{inputenc}
\usepackage{amsmath, amssymb, amsthm}
\usepackage{geometry}
\geometry{margin=2.5cm}

\newtheorem{definition}{Definición}
\newtheorem{theorem}{Teorema}
\newtheorem{lemma}{Lema}

\title{Borrador: ideas para la formalización\\de la destrucción naive de SSA}
\author{}
\date{}

\begin{document}
\maketitle

\section{Contexto}

El objetivo es demostrar que la destrucción naive de SSA que hemos implementado es correcta:
que el programa resultante (sin $\phi$-funciones) se comporta igual que el programa en forma
SSA (con $\phi$-funciones).

La implementación ya está hecha en \texttt{ssa\_destruction.rs}. El algoritmo, de forma
resumida, hace lo siguiente: para cada bloque que tiene $\phi$-funciones, y para cada
predecesor de ese bloque, mete instrucciones de copia al final del predecesor en dos fases:
\begin{enumerate}
  \item Fase 1: guarda cada valor fuente de la $\phi$ en una variable temporal nueva.
  \item Fase 2: copia cada temporal a la variable destino de la $\phi$.
\end{enumerate}
Después, borra las $\phi$-funciones del bloque.

\section{El problema central}

Las $\phi$-funciones tienen semántica \textbf{paralela}: todas leen el estado
\emph{antes} de que ninguna escriba. Es decir, si tenemos:
\[
  x = \phi(a : p_1, \ldots), \quad y = \phi(b : p_1, \ldots)
\]
ambas leen $a$ y $b$ del estado $\sigma$ original. Ninguna ve la escritura de la otra.

Pero las instrucciones de copia normales son \textbf{secuenciales}: cada una ve las
escrituras de las anteriores. Si simplemente ponemos $x \leftarrow a;\; y \leftarrow b$,
normalmente no hay problema. Pero si $b = x$ (la fuente de la segunda es el destino de
la primera), entonces la segunda leería el valor nuevo de $x$, no el original. Eso sería
incorrecto.

Por eso usamos temporales: primero guardamos todo, luego asignamos todo.

\section{Por qué funcionan los temporales}

La idea clave es que los temporales son \textbf{nombres frescos}: no existen en el
programa original. Entonces:

\begin{itemize}
  \item En la Fase 1, cuando escribimos \texttt{tmp\_0 = a}, no estamos machacando
    ninguna variable fuente (porque \texttt{tmp\_0} es un nombre nuevo, distinto de
    \texttt{a}, \texttt{b}, etc.).
  \item En la Fase 2, cuando escribimos \texttt{x = tmp\_0}, no estamos machacando
    ningún temporal (porque \texttt{x} es una variable SSA, y los temporales son
    nombres diferentes).
\end{itemize}

Resultado: cada destino $x_i$ acaba con el valor original de su fuente $y_i$, exactamente
como haría la $\phi$ en paralelo.

Esto debería poder formalizarse como un lema: algo del tipo ``la ejecución secuencial
de las $2m$ copias (Fase 1 + Fase 2) produce el mismo resultado sobre las variables
$x_1, \ldots, x_m$ que la ejecución paralela de las $\phi$-funciones''.

La demostración sería por inducción sobre el número de copias, usando que la frescura
de los temporales garantiza que las escrituras no interfieren con las lecturas posteriores.

\section{Idea de la demostración principal}

La idea es demostrar por \textbf{inducción sobre la traza} (la secuencia de bloques
que se ejecutan) que ambos programas ($G$ con phis, $G'$ sin phis) producen el mismo
resultado.

\subsection{Qué es una traza}

Una traza es simplemente la secuencia de configuraciones (estado, bloque actual, bloque
anterior) que se va generando al ejecutar el programa bloque a bloque. En cada paso se
ejecuta un bloque: primero sus $\phi$-funciones (si las tiene), luego sus instrucciones,
y luego se determina el siguiente bloque.

\subsection{El invariante}

La idea es definir un invariante $I(n)$ que diga:

\begin{quote}
  En el paso $n$ de la traza, $G$ y $G'$ están en el mismo bloque, y las variables que
  importan para ejecutar las instrucciones de ese bloque tienen los mismos valores en
  ambos.
\end{quote}

Más formalmente, al llegar al bloque $j_n$:
\begin{enumerate}
  \item Mismo camino: $j_n = j_n'$ (ambos van al mismo bloque).
  \item Mismos valores relevantes: para toda variable $v$ que aparezca como operando en
    las instrucciones del bloque, el valor de $v$ justo antes de ejecutar las instrucciones
    es igual en $G$ y en $G'$.
\end{enumerate}

La diferencia entre $G$ y $G'$ es \emph{cuándo} se asignan los valores de las $\phi$:
\begin{itemize}
  \item En $G$: las $\phi$ se ejecutan al \textbf{entrar} en el bloque destino.
  \item En $G'$: las copias equivalentes se ejecutan al \textbf{salir} del bloque predecesor.
\end{itemize}
Pero por el lema de las dos fases, el efecto es el mismo.

\subsection{Esquema del paso inductivo}

Suponiendo $I(n)$:

\begin{enumerate}
  \item Las instrucciones originales del bloque $j_n$ son idénticas en $G$ y $G'$, y por
    $I(n)$ reciben los mismos valores de entrada. Luego producen los mismos resultados.
  \item El sucesor no cambia (los saltos son los mismos), así que ambos van al mismo bloque
    siguiente. Esto da $I(n+1)$ parte (1).
  \item En $G'$, después de las instrucciones originales, se ejecutan las copias (Fase 1 + 2).
    Por el lema de las dos fases, las variables de salida de las $\phi$ del siguiente bloque
    quedan con los valores correctos. Esto da $I(n+1)$ parte (2).
\end{enumerate}

\subsection{Caso base}

$n = 0$: ambos empiezan en el bloque de entrada con los mismos valores. El bloque de entrada
no tiene predecesores, así que no tiene $\phi$-funciones. No hay diferencia. Trivial.

\section{Una propiedad útil de la construcción}

Revisando el código de construcción de SSA, he observado que un bloque con sucesor condicional
(un \texttt{if}) \textbf{nunca tiene ambos sucesores con $\phi$-funciones}. Esto es porque:

\begin{itemize}
  \item Los saltos condicionales solo se crean en \texttt{handle\_if}.
  \item Los dos sucesores inmediatos (\texttt{then} y \texttt{else}) son bloques recién creados,
    cada uno con un único predecesor.
  \item Con un solo predecesor, no se insertan $\phi$-funciones.
  \item Excepción: el \texttt{if} sin \texttt{else}, donde el sucesor falso es directamente
    el \texttt{join}. El \texttt{join} tiene 2+ predecesores (puede tener $\phi$), pero el
    \texttt{then} sigue teniendo solo 1.
\end{itemize}

Esto simplifica la demostración: no hay que preocuparse del caso en que se meten copias para
$\phi$-funciones de dos sucesores distintos en el mismo bloque predecesor. Solo hay copias
para un sucesor como máximo.

Si quisiéramos ser más generales (por si en el futuro el CFG puede venir de otra fuente),
habría que demostrar un lema adicional de ``no interferencia'' entre las copias de distintos
sucesores. Pero para nuestra construcción concreta no hace falta.

\section{Pendiente / próximos pasos}

\begin{itemize}
  \item Formalizar bien las definiciones: valuación, instrucción, bloque básico, CFG, traza.
  \item Definir formalmente la semántica paralela de las $\phi$ y la secuencial de las copias.
  \item Escribir y demostrar el lema de las dos fases (temporales emulan ejecución paralela).
  \item Escribir y demostrar el teorema principal por inducción con el invariante descrito.
  \item Definir la noción de equivalencia de trazas (mismos bloques visitados, mismas
    valuaciones proyectadas sobre variables originales).
  \item Quizás formalizar la noción de ``variable fresca'' y la relación entre nombres SSA
    y nombres originales ($\rho$).
\end{itemize}

\end{document}
